\section{侦查}
\begin{enumerate}
  \item 9 种手段和 5 种措施，每年必考其中两种。
  \item 侦查机关收集证据、查明案情、抓捕嫌疑人的相关活动。\zd{找证据，抓人}。
  \item 在有效期限内，\zd{需要变更}技术侦查措施种类或者适用对象的，\zd{应当重新办理批准手续}。
  \item \zd{事前审查 = 先批再干。事后审查 = 干了再说}。
  \item 违法侦查行为。
        \begin{enumerate}
          \item 期满不放人，不退保证金，违法查扣冻，贪挪私调用。
          \item 谁违法找谁申诉、控告 \yjt 同级 \yjt 上级检察院。
        \end{enumerate}
  \item 主体、对象、具体的适用程序。
  \item 侦查监督见“基本原则的检察院依法对刑事诉讼实行法律监督”。
  \item 讯问嫌疑人。\zd{讯问羁押仅在看守所，未押可在现住办}。\zd{核对翻译暴限胁；场外音像无法代}。
        \begin{enumerate}
          \item 送看后 \yjt 只能在看守所。没有 \yjt 现场、住处、办案场所。\al{所在市县的公司进行讯问 \cuo}。
          \item 拘传、传唤立即（目的就是为了讯问）。拘留、速捕 24 小时内。
          \item 同案犯应当分别进行。
          \item 特殊对象。选项要出现\zd{应当}。\al{未成年人、盲聋哑}。\al{讯问女性未成年，应当有女检察人员参加 \dui}。
          \item \zd{讯问时可以随便“口嗨”}。\al{把你妻子抓进来 \dui}。
          \item 嫌疑人请求\zd{自行书写供述的，应当准许}。
          \item \zd{死无重大录音 + 像；职务犯罪都录像}。
          \item 身份混同时，应当分别进行。\al{甲既是被害人，又是犯罪嫌疑人，可在讯问过程中一并收集被害人陈述 \cuo}。
        \end{enumerate}
  \item 询问证人、被害人。\zd{询证提住侦现单}。\zd{核对翻译暴限胁；毒醉猜测未个别；拒绝出庭 + 不真实}。
        \begin{enumerate}
          \item 证人提出的地点、住处；侦查现场；单位。\zd{不能指定地点}。
          \item 没有规定要什么时候开始询问证人。
          \item 没有规定一次讯问可以持续多久，但是一般认为最长不得超过 12 小时。
          \item 个别进行。
          \item \zd{有社会影响的案件应当全程录音录像}。
          \item 特殊对象。选项要出现\zd{应当}。\al{未成年人、盲聋哑}。
          \item \zd{讯问时不得口嗨}。
        \end{enumerate}
  \item 勘验（死的）、检查（活的）。\zd{侦查人员}对犯罪有关场所、物品、尸体、人身进行勘查或检验。\zd{无法解释}。
        \begin{enumerate}
          \item 检查。\zx{可以指派、聘请具有专门知识的人}在侦查人员\zd{主持下}进行。
          \item \zd{勘察证没有例外}。
          \item 现场勘验、物证检验、尸体检验、人身检查、强制取样、侦查实验。
          \item 女性由\zd{女性工作人员或医师}。\al{应当由女性工作人员或女医师进行 \cuo{不要求性别}}。
          \item \zd{尸体解剖通知即可，无需同意}。
          \item \zd{必要时可以进行侦察实验，进行侦察实验，应当全程录音录像}。
        \end{enumerate}
  \item 搜查。\zx{一定是侦查人员}对嫌疑人人身、物品、住处或者有关的场所进行搜查、检查。\zd{无法解释}。
        \begin{enumerate}
          \item \zd{女性只能由女性工作人员}。
          \item 紧急情况，\zd{可以没有搜查证}。
        \end{enumerate}
  \item 查扣冻。扣押物品清单，一式两份。\zd{喜欢结合没收程序、涉案财物处理一起考}。
        \begin{enumerate}
          \item 邮件电报（通信自由）要经公安 \hz 检察院批准。\al{可以通知 \cuo}。
          \item 本案无关，\zx{3 日解除}。
        \end{enumerate}
  \item 鉴定。\zd{鉴定人负责制度}。\zd{见错就排}（可替代）。
        \begin{enumerate}
          \item 指派 \hz 聘请（不能委托）、自然人、书面、独立进行、签名 \hz 盖章、\zd{可自费要求补充 \hz 重新鉴定}。
          \item \zd{精神病鉴定时间不计入办案期限}。
          \item 法官决定是否采纳。
        \end{enumerate}
  \item 辨认。\zd{不是侦调来主持，指示预见个混混}。
        \begin{enumerate}
          \item 主体。被害人、证人、嫌疑人。要个别。
          \item 对象。物品、文件、场所。\zd{证人和被害人不能成为辨认对象}。要混杂。
          \item （主体）单独原则。（对象）混杂（同类）原则。总数。\zd{公 7 人 5 物 10 照。检 7 人 5 物 10 人照 5 物照}。
          \item \zd{防止预断}。不能提前见，但要提前询问被辨认对象的特征。
          \item 禁止暗示、现场监督、制作笔录、保密原则。
          \item 能够准确描述\zd{独有特征}的，\zd{陪衬物数量不受限制}。\al{场所、尸体}。
        \end{enumerate}
  \item 技术侦查。电子侦听；电话监听；电子监控；秘密拍照 \hz 录像；秘密获取某些物证、邮件。
        \begin{enumerate}
          \item \zd{公安危恐黑毒严，检院用权侵人权，监委贪贿职务犯，都有追捕在逃犯，非常重大、危险、恐怖}。
          \item 应当在\zd{立案后}实施。公检决定\zd{公安执行}！！\al{只有在立案后，才能 xx \dui}。
                \zd{用一长串正确的表述麻痹你，然后大海捞针}。
                \al{检可以利用对职权实施的严重侵犯公民人身权利的重大犯罪案件\xhx{实施}技术侦查 \cuo}。
          \item 措施报告书 \yjt 设区的市一级以上公安负责人批准 \yjt 措施决定书。\al{需要经过省公安批准 \cuo}。
          \item 有效期 3 + n 个 3。\al{不得超过 9 个月 \cuo}。
          \item 秘密措施。\zd{卧底侦查、隐匿身份（不能引诱犯罪）、控制下交付、诱惑侦查、跟踪}。
                \al{批准采取了监听，实施跟踪措施无需再报请审批 \cuo}。
          \item \zd{区分启动技侦和启动后秘密措施决定的主体。启动后负责人批即可}。
          \item \al{需要经过转化才能使用、涉及国家秘密不能采用 \cuo。可做定案依据，但要调查 \hz 庭外予以核实 \dui}。
          \item 得来的证据\zd{保密 \hz 销毁。单说一个都是错的}。
        \end{enumerate}
  \item 通缉。公安通令缉拿\zd{应当逮捕却在逃的}犯罪嫌疑人。
        \begin{enumerate}
          \item 公检监决定，但\zd{只能公安发布}通缉令。
          \item \zd{我的地盘听我的，超出辖区找上级；全国各找各妈}。
        \end{enumerate}
  \item 限制出境。省以上监委批准，公安执行。
  \item \zd{扣押查验认}。检查、搜查、查封；勘验；辨认；扣押。必须要有见证人。\zd{有缺陷，有利害，有职权就排除}。
        \al{甲称曾举报乙贪污，乙报案是对他的打击报复。可由乙在场担任见证人 \cuo}。
  \item 可以没有证的。\zx{搜、拘、唤}。搜查、拘留、先行拘留、传唤（唯一可口头）。
\end{enumerate}
